\chapter{Luces fuera}
Este es el juego de luces fuera de la aplicación.
Es un tablero de 5 por 5 y el objetivo es apagar todas las casillas.
Al pulsar una, cambian ella y las que tiene arriba, abajo, a la izquierda y a la derecha.

Se sigue la misma paleta de colores que en el resto de la aplicación.

\begin{figure}[H]
  \centering
  \includegraphics[width=0.4\textwidth]{luces}
\end{figure}

\section{Representación del tablero}
El estado del juego se guarda en una matriz de números, donde $0$ es apagada y $1$ es encendida.

\begin{lstlisting}
private val luces = Array(tamano) { IntArray(tamano) }
\end{lstlisting}

Otra matriz guarda los \texttt{TextView} en el mismo orden, así que la misma fila y columna sirven para las dos.

\section{Nueva partida}
La función \texttt{nuevaPartida} apaga todo el tablero y después simula 5 pulsaciones al azar.

\begin{lstlisting}
for (i in 1..5) {
    cambiarCruz(Random.nextInt(tamano), Random.nextInt(tamano))
}
\end{lstlisting}

Así el tablero siempre tiene solución, porque basta con repetir esas pulsaciones.
Si se pusieran las luces al azar, muchas veces no se podría resolver.

\section{Cambiar las luces}
La función \texttt{cambiarCruz} llama a \texttt{cambiar} para la casilla pulsada y las cuatro de alrededor.

\begin{lstlisting}
private fun cambiar(fila: Int, columna: Int) {
    if (fila in 0 until tamano && columna in 0 until tamano) {
        if (luces[fila][columna] == 0) {
            luces[fila][columna] = 1
        } else {
            luces[fila][columna] = 0
        }
    }
}
\end{lstlisting}

Antes de cambiar una casilla, \texttt{cambiar} comprueba que está dentro del tablero.
De esta forma los bordes y las esquinas no dan error.

\section{Pulsar una casilla y ganar}
Al tocar una casilla se llama a \texttt{pulsar}, que cambia las luces y actualiza la pantalla.
Después comprueba si se ha ganado con \texttt{todasApagadas}.
Esta función recorre la matriz y devuelve \texttt{false} en cuanto encuentra una luz encendida.
Si no hay ninguna, se muestra un \texttt{Toast} avisando de que se ha ganado.

\section{Actualizar la pantalla}
La función \texttt{actualizarPantalla} copia la matriz a las casillas.
Pone el número y los colores de cada una: mostaza si está encendida y verde si está apagada.